% 科研示意图：白底、深灰细线与柔和填充；与封面和正文调色板分开维护。
\definecolor{FigureInk}{HTML}{222222}
\definecolor{FigureStroke}{HTML}{444444}
\definecolor{FigureMuted}{HTML}{686868}
\definecolor{FigureGrid}{HTML}{C4C9CC}
\definecolor{FigurePaper}{HTML}{F4F5F6}
% 离散强调色按种类数选择：单色黄、双色蓝红、三色蓝红黄。
\definecolor{FigureBlue}{HTML}{7998AD}
\definecolor{FigureRed}{HTML}{D57B70}
\definecolor{FigureYellow}{HTML}{D6B35D}
\definecolor{FigureBlueFill}{HTML}{EDF2F5}
\definecolor{FigureRedFill}{HTML}{F8E8E5}
\definecolor{FigureYellowFill}{HTML}{FFF1CF}
\definecolor{FigureCream}{HTML}{FBEFDA}
% 其他卷保留原有默认色；第二卷在局部分组中覆盖。
\colorlet{FigureInput}{BookBlue}
\colorlet{FigureModel}{BookTeal}
\colorlet{FigureOutput}{BookInk}
\colorlet{FigureLoss}{BookAmber}
% 神经网络图：白色画布与浅色节点保留拓扑，明亮强调色和清晰线条标出计算职责。
\definecolor{NNInput}{HTML}{2687EA}
\definecolor{NNHidden}{HTML}{9464DA}
\definecolor{NNOutput}{HTML}{12ADA7}
\definecolor{NNGradient}{HTML}{A95B13}
\definecolor{NNLine}{HTML}{8A9CB6}
\colorlet{NNInk}{BookInk}
\colorlet{NNCanvas}{white}
\colorlet{NNPanel}{NNInput!4!white}
\definecolor{NNWire}{HTML}{657B99}
\colorlet{FigureInputFill}{FigureInput!22!white}
\colorlet{FigureModelFill}{FigureModel!18!white}
\colorlet{FigureOutputFill}{FigureOutput!20!white}
\colorlet{FigureLossFill}{FigureLoss!18!white}
\colorlet{FigurePanel}{FigureModel!9!white}
\colorlet{FoundationInput}{FigureInput}
\colorlet{FoundationModel}{FigureModel}
\colorlet{FoundationOutput}{FigureOutput}
\colorlet{FoundationLoss}{FigureLoss}
